
\documentclass{article}
\usepackage{colortbl}
\usepackage{makecell}
\usepackage{multirow}
\usepackage{supertabular}

\begin{document}

\newcounter{utterance}

\centering \large Interaction Transcript for game `cladder', experiment `full\_v1.5\_default', episode 5930 with reflexion{-}llama{-}3.1{-}8b{-}instruct.
\vspace{24pt}

{ \footnotesize  \setcounter{utterance}{1}
\setlength{\tabcolsep}{0pt}
\begin{supertabular}{c@{$\;$}|p{.15\linewidth}@{}p{.15\linewidth}p{.15\linewidth}p{.15\linewidth}p{.15\linewidth}p{.15\linewidth}}
    \# & \multicolumn{2}{c}{Player} && \multicolumn{2}{c}{Game Master} \\
    \hline

    \theutterance \stepcounter{utterance}  
    & & & \multicolumn{4}{p{0.6\linewidth}}{
        \cellcolor[rgb]{0.9,0.9,0.9}{
            \makecell[{{p{\linewidth}}}]{
                \texttt{\tiny{[P1$\langle$GM]}}
                \texttt{You have attempted this same task before and did not succeed. These are the plans you wrote after those attempts.} \\
\texttt{Trial \#1: New plan: } \\
\\ 
\texttt{In my previous attempt, I correctly identified the causal chain and the mediating effect of the wife's action. However, I failed to consider the possibility that the wife's action might be influenced by other factors, such as her own habits or preferences, that are not mentioned in the conditions. } \\
\\ 
\texttt{To improve my strategy, I will consider the possibility of additional, unmentioned factors that might influence the wife's action. I will also pay closer attention to the wording of the conditions, looking for any subtle implications or nuances that might affect the causal relationships. } \\
\\ 
\texttt{Specifically, I will ask myself questions like: Are there any other factors that might influence the wife's action of setting the alarm? Are there any indirect effects or secondary relationships that might affect the outcome? By considering these additional factors and nuances, I hope to develop a more comprehensive and accurate understanding of the causal relationships in the hypothetical world.} \\
\\ 
\\ 
\texttt{You are an expert in causal inference. The following question is not a typical commonsense query, but rather a meticulously designed question created by a professor specializing in causal inference, intended to assess the students' mastery of the course content.} \\
\texttt{Imagine a self{-}contained, hypothetical world with only the following conditions, and without any unmentioned factors or causal relationships: Drinking coffee has a direct effect on wife and alarm clock. Wife has a direct effect on alarm clock. We know that drinking coffee causes alarm set by wife. drinking coffee or alarm set by wife causes ringing alarm. Would the alarm rings the next morning if drinking coffee instead of not drinking coffee?} \\
\texttt{Start your answer with 'Yes' or 'No', followed by additional reasoning or evidence to support your explanation.} \\
            }
        }
    }
     \\ \\

    \theutterance \stepcounter{utterance}  
    & \multicolumn{4}{p{0.6\linewidth}}{
        \cellcolor[rgb]{0.561,0.737,0.561}{
            \makecell[{{p{\linewidth}}}]{
                \texttt{\tiny{[P1$\rangle$GM]}}
                \texttt{Yes.} \\
\\ 
\texttt{The given conditions establish a causal chain from drinking coffee to alarm set by wife, and from alarm set by wife to ringing alarm. The direct effect of drinking coffee on wife and alarm clock is not specified to be relevant to the outcome of the alarm ringing. } \\
\\ 
\texttt{The condition "drinking coffee causes alarm set by wife" implies that drinking coffee is a sufficient cause for the wife to set the alarm. The condition "drinking coffee or alarm set by wife causes ringing alarm" implies that either drinking coffee or the wife setting the alarm is a sufficient cause for the alarm to ring. } \\
\\ 
\texttt{Since drinking coffee is a sufficient cause for the wife to set the alarm, and the wife setting the alarm is a sufficient cause for the alarm to ring, it follows that drinking coffee is a sufficient cause for the alarm to ring. Therefore, if drinking coffee instead of not drinking coffee, the alarm will ring the next morning.} \\
            }
        }
    }
    & & \\ \\

    \theutterance \stepcounter{utterance}  
    & & & \multicolumn{2}{p{0.3\linewidth}}{
        \cellcolor[rgb]{0.9,0.9,0.9}{
            \makecell[{{p{\linewidth}}}]{
                \texttt{\tiny{[GM$|$GM]}}
                \texttt{Yes} \\
            }
        }
    }
    & & \\ \\

    \theutterance \stepcounter{utterance}  
    & & & \multicolumn{2}{p{0.3\linewidth}}{
        \cellcolor[rgb]{0.9,0.9,0.9}{
            \makecell[{{p{\linewidth}}}]{
                \texttt{\tiny{[GM$|$GM]}}
                \texttt{yes} \\
            }
        }
    }
    & & \\ \\

    \theutterance \stepcounter{utterance}  
    & & & \multicolumn{2}{p{0.3\linewidth}}{
        \cellcolor[rgb]{0.9,0.9,0.9}{
            \makecell[{{p{\linewidth}}}]{
                \texttt{\tiny{[GM$|$GM]}}
                \texttt{game\_result = WIN} \\
            }
        }
    }
    & & \\ \\

\end{supertabular}
}

\end{document}
